% GENERATED FILE. Edit canonical lessons, then run npm run generate:books.
% canonical-source-hash: fnv1a64:44b9ff2207317088
% canonical-lessons: TE-C496-badhapadu, TE-C496-badhapettu, TE-C496-kungipovu, TE-C496-kumilipovu, TE-C496-nisprha

\chapter{Sorrow}
\label{ch:te-sorrow}
\hlchaptermodality{496}

\begin{chapteropening}
\textbf{By the end of this chapter:} \emph{I can say that I was sad, hurt or crushed.}

Complete the last lesson of chapter 496: I can say that I was sad, hurt or crushed.
\end{chapteropening}

\section[\texorpdfstring{\te{బాధపడు} (bādhapaḍu)}{బాధపడు (bādhapaḍu)}]{\te{బాధపడు} (bādhapaḍu) — to feel sad; to grieve}

\label{lesson:TE-C496-badhapadu}



\begin{quote}
Before the new one: say the Telugu for a nightmare, then the Telugu for to be distraught.
\end{quote}

\subsection*{You'll want to know: \te{బాధపడు}}
\textbf{\te{బాధపడు}} — \emph{bādhapaḍu} — \textquotedblleft{}to feel sad; to grieve\textquotedblright{}.

\begin{quote}
\te{స్నేహితుడు} \te{వెళ్ళిపోయాక} \te{నేను} \te{చాలా} \te{బాధపడ్డాను}.
\end{quote}

\begin{quote}
\emph{snēhituḍu veḷḷipōyāka nēnu cālā bādhapaḍḍānu.}
\end{quote}

\begin{quote}
After my friend left, I was very sad.
\end{quote}

Past (I …): \te{బాధపడ్డాను}, \emph{bādhapaḍḍānu}.

The cause takes -\te{కి}: \te{ఆ} \te{మాటకి} \te{బాధపడ్డాను}, that remark upset me.

\begin{cousinweb}
\textbf{\te{బాధపడు}} (\emph{bādhapaḍu}) joins \te{బాధ}, \textquotedblleft{}pain, sorrow\textquotedblright{}, a loan from Sanskrit \emph{bādhā}, to the native verb \emph{paḍu}, \textquotedblleft{}to feel, to suffer\textquotedblright{} (DEDR 3853).
\end{cousinweb}

\begin{grammarlens}[title={What you've built}]
One more word for describing what you lived and felt.
\end{grammarlens}

\subsection*{Your turn}
\begin{itemize}
\raggedright
  \item \emph{Say it:} \emph{bādhapaḍu}
  \item \emph{Say it:} \emph{bādhapaḍu}, once more
  \item \emph{Say it:} \emph{snēhituḍu veḷḷipōyāka nēnu cālā bādhapaḍḍānu.} — After my friend left, I was very sad.
  \item \emph{Recall:} say the Telugu for a nightmare, then the Telugu for to be distraught, then say \emph{bādhapaḍu} again
\end{itemize}

\begin{tcolorbox}[breakable,colback=teal!4,colframe=teal!35!black,title={Before you move on}]
What does \te{బాధపడు} mean? (\textquotedblleft{}To feel sad; to grieve\textquotedblright{}.) Say it once more.
\end{tcolorbox}
\section[\texorpdfstring{\te{బాధపెట్టు} (bādhapeṭṭu)}{బాధపెట్టు (bādhapeṭṭu)}]{\te{బాధపెట్టు} (bādhapeṭṭu) — to hurt; to upset}

\label{lesson:TE-C496-badhapettu}



\begin{quote}
Before the new one: say the Telugu for to be distraught, then the Telugu for to feel sad.
\end{quote}

\subsection*{You'll want to know: \te{బాధపెట్టు}}
\textbf{\te{బాధపెట్టు}} — \emph{bādhapeṭṭu} — \textquotedblleft{}to hurt; to upset\textquotedblright{}.

\begin{quote}
\te{నా} \te{మాటలు} \te{ఆమెను} \te{బాధపెట్టాయి}.
\end{quote}

\begin{quote}
\emph{nā māṭalu āmenu bādhapeṭṭāyi.}
\end{quote}

\begin{quote}
My words hurt her.
\end{quote}

Past (I …): \te{బాధపెట్టాను}, \emph{bādhapeṭṭānu}.

For feelings rather than the body: \te{మనసును} \te{బాధపెట్టు}, to hurt someone deeply.

\begin{cousinweb}
\textbf{\te{బాధపెట్టు}} (\emph{bādhapeṭṭu}) joins \te{బాధ}, \textquotedblleft{}pain\textquotedblright{}, a loan from Sanskrit \emph{bādhā}, to the native verb \emph{peṭṭu}, \textquotedblleft{}to put\textquotedblright{}: to cause someone pain.
\end{cousinweb}

\begin{grammarlens}[title={What you've built}]
One more word for describing what you lived and felt.
\end{grammarlens}

\subsection*{Your turn}
\begin{itemize}
\raggedright
  \item \emph{Say it:} \emph{bādhapeṭṭu}
  \item \emph{Say it:} \emph{bādhapeṭṭu}, once more
  \item \emph{Say it:} \emph{nā māṭalu āmenu bādhapeṭṭāyi.} — My words hurt her.
  \item \emph{Recall:} say the Telugu for to be distraught, then the Telugu for to feel sad, then say \emph{bādhapeṭṭu} again
\end{itemize}

\begin{tcolorbox}[breakable,colback=teal!4,colframe=teal!35!black,title={Before you move on}]
What does \te{బాధపెట్టు} mean? (\textquotedblleft{}To hurt; to upset\textquotedblright{}.) Say it once more.
\end{tcolorbox}
\section[\texorpdfstring{\te{కుంగిపోవు} (kuṅgipōvu)}{కుంగిపోవు (kuṅgipōvu)}]{\te{కుంగిపోవు} (kuṅgipōvu) — to be crushed; to sink in spirit}

\label{lesson:TE-C496-kungipovu}



\begin{quote}
Before the new one: say the Telugu for to feel sad, then the Telugu for to hurt.
\end{quote}

\subsection*{You'll want to know: \te{కుంగిపోవు}}
\textbf{\te{కుంగిపోవు}} — \emph{kuṅgipōvu} — \textquotedblleft{}to be crushed; to sink in spirit\textquotedblright{}.

\begin{quote}
\te{ఉద్యోగం} \te{పోయాక} \te{అతను} \te{కుంగిపోయాడు}.
\end{quote}

\begin{quote}
\emph{udyōgaṁ pōyāka atanu kuṅgipōyāḍu.}
\end{quote}

\begin{quote}
After losing his job, he was crushed.
\end{quote}

Past (I …): \te{కుంగిపోయాను}, \emph{kuṅgipōyānu}.

Literally to sink and go: a -\te{పోవు} compound, like \te{వెళ్ళిపోవు}.

\begin{cousinweb}
\textbf{\te{కుంగిపోవు}} (\emph{kuṅgipōvu}) is a native compound: DEDR 1767 lists Telugu \emph{k(r)uṅgu}, \textquotedblleft{}to sink, to fall down\textquotedblright{}, with Kannada \emph{kuggu}; \emph{pōvu}, \textquotedblleft{}to go\textquotedblright{} (DEDR 4572) completes it. Someone who sinks in spirit is crushed.
\end{cousinweb}

\begin{grammarlens}[title={What you've built}]
One more word for describing what you lived and felt.
\end{grammarlens}

\subsection*{Your turn}
\begin{itemize}
\raggedright
  \item \emph{Say it:} \emph{kuṅgipōvu}
  \item \emph{Say it:} \emph{kuṅgipōvu}, once more
  \item \emph{Say it:} \emph{udyōgaṁ pōyāka atanu kuṅgipōyāḍu.} — After losing his job, he was crushed.
  \item \emph{Recall:} say the Telugu for to feel sad, then the Telugu for to hurt, then say \emph{kuṅgipōvu} again
\end{itemize}

\begin{tcolorbox}[breakable,colback=teal!4,colframe=teal!35!black,title={Before you move on}]
What does \te{కుంగిపోవు} mean? (\textquotedblleft{}To be crushed; to sink in spirit\textquotedblright{}.) Say it once more.
\end{tcolorbox}
\section[\texorpdfstring{\te{కుమిలిపోవు} (kumilipōvu)}{కుమిలిపోవు (kumilipōvu)}]{\te{కుమిలిపోవు} (kumilipōvu) — to grieve inwardly}

\label{lesson:TE-C496-kumilipovu}



\begin{quote}
Before the new one: say the Telugu for to hurt, then the Telugu for to be crushed.
\end{quote}

\subsection*{You'll want to know: \te{కుమిలిపోవు}}
\textbf{\te{కుమిలిపోవు}} — \emph{kumilipōvu} — \textquotedblleft{}to grieve inwardly\textquotedblright{}.

\begin{quote}
\te{ఆమె} \te{ఎవరికీ} \te{చెప్పకుండా} \te{కుమిలిపోయింది}.
\end{quote}

\begin{quote}
\emph{āme evarikī ceppakuṇḍā kumilipōyindi.}
\end{quote}

\begin{quote}
She grieved inside without telling anyone.
\end{quote}

Past (I …): \te{కుమిలిపోయాను}, \emph{kumilipōyānu}.

Grief that smoulders without being spoken.

\begin{cousinweb}
\textbf{\te{కుమిలిపోవు}} (\emph{kumilipōvu}) is native Dravidian: DEDR 1752 lists Telugu \emph{kumulu}, \textquotedblleft{}to smoulder without flame, to be consumed inwardly, to grieve\textquotedblright{}, and the verb ends in \emph{pōvu}, \textquotedblleft{}to go\textquotedblright{} (DEDR 4572).
\end{cousinweb}

\begin{grammarlens}[title={What you've built}]
One more word for describing what you lived and felt.
\end{grammarlens}

\subsection*{Your turn}
\begin{itemize}
\raggedright
  \item \emph{Say it:} \emph{kumilipōvu}
  \item \emph{Say it:} \emph{kumilipōvu}, once more
  \item \emph{Say it:} \emph{āme evarikī ceppakuṇḍā kumilipōyindi.} — She grieved inside without telling anyone.
  \item \emph{Recall:} say the Telugu for to hurt, then the Telugu for to be crushed, then say \emph{kumilipōvu} again
\end{itemize}

\begin{tcolorbox}[breakable,colback=teal!4,colframe=teal!35!black,title={Before you move on}]
What does \te{కుమిలిపోవు} mean? (\textquotedblleft{}To grieve inwardly\textquotedblright{}.) Say it once more.
\end{tcolorbox}
\section[\texorpdfstring{\te{నిస్పృహ} (nispṛha)}{నిస్పృహ (nispṛha)}]{\te{నిస్పృహ} (nispṛha) — despair; dejection}

\label{lesson:TE-C496-nisprha}



\begin{quote}
Before the new one: say the Telugu for to be crushed, then the Telugu for to grieve inwardly.
\end{quote}

\subsection*{You'll want to know: \te{నిస్పృహ}}
\textbf{\te{నిస్పృహ}} — \emph{nispṛha} — \textquotedblleft{}despair; dejection\textquotedblright{}.

\begin{quote}
\te{చాలా} \te{ప్రయత్నాల} \te{తరువాత} \te{అతనికి} \te{నిస్పృహ} \te{కలిగింది}.
\end{quote}

\begin{quote}
\emph{cālā prayatnāla taruvāta ataniki nispṛha kaligindi.}
\end{quote}

\begin{quote}
After many attempts, despair set in.
\end{quote}

A formal word: \te{నిస్పృహ} \te{కలుగు}, despair sets in.

\begin{cousinweb}
\textbf{\te{నిస్పృహ}} (\emph{nispṛha}) is borrowed from Sanskrit \emph{niḥspṛha}, \textquotedblleft{}free from desire, not longing for anything\textquotedblright{} (Monier-Williams), from \emph{niḥ-}, \textquotedblleft{}without\textquotedblright{}, and \emph{spṛhā}, \textquotedblleft{}longing\textquotedblright{}. In Telugu, having nothing left to long for is despair.
\end{cousinweb}

\begin{grammarlens}[title={What you've built}]
One more word for describing what you lived and felt.
\end{grammarlens}

\subsection*{Your turn}
\begin{itemize}
\raggedright
  \item \emph{Say it:} \emph{nispṛha}
  \item \emph{Say it:} \emph{nispṛha}, once more
  \item \emph{Say it:} \emph{cālā prayatnāla taruvāta ataniki nispṛha kaligindi.} — After many attempts, despair set in.
  \item \emph{Recall:} say the Telugu for to be crushed, then the Telugu for to grieve inwardly, then say \emph{nispṛha} again
\end{itemize}

\begin{tcolorbox}[breakable,colback=teal!4,colframe=teal!35!black,title={Before you move on}]
What does \te{నిస్పృహ} mean? (\textquotedblleft{}Despair; dejection\textquotedblright{}.) Say it once more.
\end{tcolorbox}
